% SPDX-License-Identifier: GPL-3.0-or-later
% Japanese -- what the reading editions' preamble leaves to the language.
% Input by lib/frank-preamble.tex as \input{\FrankLib/lang/\BookLang.tex};
% docs/languages.md section 6 lists the hooks every file here must define.
% The CJK model: a reading (kana) over every word -- or over every chunk, for
% a chunk written without its words -- the reading alone as a pass of its own,
% no marks to strip, and a last pass that is not a font but a direction -- the
% text set vertically.
%
% There is no luatexja on this machine, and the editions are built with
% LuaLaTeX, so everything Japanese here is done with babel and fontspec alone:
% babel's `japanese` locale with intraspace glue for the line breaks (Japanese
% has no spaces to break at), and the vertical pass as a rotated box in a
% font loaded with the OpenType `vert` feature -- the same trick the studio
% uses under XeLaTeX with Vertical=RotatedGlyphs (docs/languages.md section 5).
% Both were tried on this machine before this file was written.

% ---- the babel language and the switches ----------------------------------
\newcommand{\FrankLangName}{japanese}
\FrankRTLfalse           % chunk left, gloss right; \pw needs no bidi isolate
\FrankHasBaretrue        % pass 4 is the same text without the furigana
\FrankHasAlttrue         % an alternate face exists (gothic): chapter numbers, the title
\FrankHasReadingtrue     % \chr carries the kana: over the chunk in pass 1, heading the gloss
\FrankHasVerttrue        % the last pass is \FrankVertPass, not the font pass
\FrankHasAloudtrue       % pass 2 is the kana alone, chunk after chunk

% ---- fonts ------------------------------------------------------------------
% From the registry's "tex" record: Noto Serif CJK JP (the system's; nothing
% CJK travels with the toolbox, the files are too large), then the fallbacks.
% intraspace=0 .1 0 is the glue babel puts between two CJK characters: no
% natural width, a tenth of an em of stretch, no shrink -- what lets a line
% break anywhere and still be justified.  Renderer=HarfBuzz as for the Arabic
% script, so one engine shapes every language.  Script=CJK is fontspec's name
% for the hani/kana scripts.
\babelprovide[import=ja,onchar=ids fonts,intraspace=0 .1 0]{japanese}
\FrankPickFont{\FrankMainFontName}{ja}{main}{main_fallbacks}
\babelfont[japanese]{rm}[Script=CJK,Language=Japanese,Renderer=HarfBuzz]{\FrankMainFontName}
% The alternate face is the gothic (sans).  The PDF's fourth pass is the
% vertical one, so this face is only used where \ifFrankHasAlt asks for it:
% the chapter number and the title page.  The HTML reader labels its own
% fourth pass with it (docs/languages.md section 6).
\FrankPickFont{\FrankAltFontName}{ja}{alt}{alt_fallbacks}
\newfontfamily\FrankGothicFont[Script=CJK,Language=Japanese,Renderer=HarfBuzz]{\FrankAltFontName}
\newcommand{\FrankAltFont}{\FrankGothicFont}
% The vertical face: the main font again, with the `vert` and `vrt2` features
% on, so that the glyphs which change shape in vertical text -- the brackets,
% the long-vowel mark, the small kana, the comma and the full stop -- are the
% vertical forms.  Set horizontally in this font and turned clockwise, each
% glyph then stands upright in its column.  HarfBuzz honours the features on
% this machine (checked: the rotated brackets and the vertical comma appear);
% Renderer=Node does too, and additionally maps the punctuation to its
% presentation forms in the text layer, which is why HarfBuzz was kept --
% pdftotext then gives back the characters the book was written with.
\newfontfamily\FrankVertFont[Script=CJK,Language=Japanese,Renderer=HarfBuzz,%
  RawFeature={+vert,+vrt2}]{\FrankMainFontName}
\usepackage{graphicx}   % \rotatebox and \resizebox for the vertical pass

% ---- the two Lua filters ---------------------------------------------------
% Identity: the kana lives in its own field, so the bare text is the text
% itself (the registry's `strip` is null), and the labels are in Latin digits.
\directlua{
  function frank_strip(s) return s end
  function frank_latin(s) return s end
}

% ---- the seam between two chunks --------------------------------------------
% The core joins the chunks of a sentence with \FrankChunkSep when it rebuilds
% pass 1 and the plain and vertical passes from the buffers.  A space, the default, is right for a
% language written with spaces; Japanese is not, and a space there would show
% as a gap the source does not have.  So the seam is glue of no width with the
% same stretch babel puts between two characters: invisible, and a place to
% break the line -- which the ruby boxes of pass 1 need, because babel never
% inserts its glue between two boxes, only between two characters.
\newcommand{\FrankChunkSep}{\hskip 0pt plus .1em\relax}

% ---- the vocabulary line -----------------------------------------------------
% \vb{dictionary form}{romaji}{-masu stem}{romaji}{-te form}{romaji}{meaning}:
% the labels the core prints before the two extra forms.  The stem is what
% -masu, -tai and -nagara attach to, the -te form what te, ta, tara and teiru
% are made of and the one form nobody can guess (書いて, 泳いで, 読んで, and
% 行って against every other -ku verb); the past is the -te form with ta and is
% not printed.  A -suru verb is the whole of it in every slot: 勉強する,
% 勉強し, 勉強して.  The sounds are rōmaji; the kana line of the chunk already
% carries the reading, so no kana goes into a \vb.
% What the three forms do not name goes in ONE parenthesis after the meaning,
% items parted by "; ": first the class, which is what builds every form that
% is NOT printed (買う makes 買わない, 食べる 食べない) --
%   godan | ichidan | suru | irregular      (する itself, 来る and its compounds;
%                                            a noun + する, 気にする, is suru)
% then tr. / intr. / tr./intr. where the dictionary says (開く and 開ける share
% one English gloss and not one particle), then hon. or hum. for a verb that
% is honorific or humble:
%   \vb{書く}{kaku}{書き}{kaki}{書いて}{kaite}{to write (godan; tr.)}
% lib/verbs/ja.py writes this entry for the editor's sidebar from dict/ja.db:
% the forms and their rōmaji from the head line of the Wiktionary page (the
% -te form is the past with た made て), the class from its head template,
% transitivity where the page gives one -- about half the verbs.
% The reader takes the same pair of labels from the registry
% (lib/languages.json, "vb_labels": ["stem", "-te"], with the three forms in
% words as "vb_forms"); the two must agree, or the PDF and the reader of one
% book label the same form differently.  The registry holds the word alone --
% the italics below are this page's typography, not part of the label -- and
% lib/newlang.py --check compares them that way.
\newcommand{\FrankVbPres}{stem}
\newcommand{\FrankVbPast}{\textit{-te}}

% ---- group ruby --------------------------------------------------------------
% \jruby{text}{kana}: the reading, small, centred over the chunk.  A plain
% \vbox with the text's baseline as its own; the pair is one unbreakable box,
% which is why the seam glue above matters.  The kana inherits the colour, so a
% highlighted chunk (\Cred etc.) is highlighted with its reading.  In a
% language without a reading the core's \providecommand prints the text alone.
% It is the core's \FrankRuby too: a \chrw sets each WORD's reading with this
% same box, measured the same way, and the seam between two words is glue for
% the reason the seam between two chunks is.
%
% A box cannot be broken, and the chunk column is a third of the measure: a
% chunk of nine characters at 13pt is wider than it, and its ruby box ran into
% the gloss (14 overfull boxes in the fixture edition).  So the pair is
% measured first, in vertical mode -- \setbox starts no paragraph, whereas a
% \leavevmode before the test would leave an empty line behind -- and when it
% is wider than the current measure (\hsize: the column, or the whole line in
% pass 1) it is set stacked instead: the reading on a line of its own, then
% the text as an ordinary paragraph, which babel breaks between characters.
\newcommand{\jruby}[2]{%
  \begingroup
    \setbox0\hbox{{\SzTrans #2}}%
    \setbox2\hbox{#1}%
    \dimen0=\wd0 \ifdim\wd2>\dimen0 \dimen0=\wd2 \fi
    \ifdim\dimen0>\hsize
      \par\noindent{\SzTrans #2\par}\nobreak\noindent #1%
    \else
      \leavevmode
      \vbox{\hbox to\dimen0{\hss\box0\hss}%
            \nointerlineskip\kern0.5pt
            \hbox to\dimen0{\hss\box2\hss}}%
    \fi
  \endgroup}

% ---- the vertical pass -------------------------------------------------------
% The bare text (\FrankBuf, the same buffer the plain pass sets) as tategaki: set in a
% \parbox whose WIDTH is the column height, in the vertical-forms font, then
% turned clockwise with \rotatebox.  Line 1 becomes the rightmost column and
% the columns run right to left, which is the order a Japanese book is read in.
%
% HOW TALL THE COLUMNS ARE.  The turned box cannot be broken across pages, so
% its height decides how the pages fill.  A first version set every sentence
% in columns of 0.78\textheight whatever its length: a box that tall follows
% pass 3 on the same page almost never, so every fourth pass opened a fresh
% page and left the one before it mostly white, and a sentence of a column and
% a bit left its second column mostly empty (the fixture edition came out at
% 22 pages).  So the sentence is measured first -- its natural width in the
% vertical face, \wd1, is its length as a single column -- and spread over as
% many columns as it needs at \FrankVertColumn, each column getting an equal
% share of the text plus one em of slack so that the whole-character line
% breaks can always fill the last column from the others.  A short sentence
% is then a short box that sits under pass 3; a long one has many columns of
% a comfortable height rather than few of a great one.  When that would want
% more columns than the measure holds (a column is one leading wide once
% turned) the columns grow towards \FrankVertHeight instead, and only a
% sentence too long even for that is scaled down to \linewidth, as before.
% The column count is worked out in Lua because it is a ceiling and a floor
% of quotients of dimensions, which TeX's own arithmetic (\numexpr rounds)
% makes needlessly awkward: frank_vert_columns(width of the sentence, usable
% column height, usable tallest height, measure, leading) in sp.
% \raise removes the depth the rotation gives the box, so it sits on the page
% like a paragraph and the page breaker sees its whole height.
\newcommand{\FrankVertHeight}{0.78\textheight}   % the tallest a column may ever be
\newcommand{\FrankVertColumn}{0.5\textheight}    % the height the columns are balanced to
\directlua{
  function frank_vert_columns(w, column, tallest, measure, leading)
    local most = math.floor(measure / leading)
    local n = math.ceil(w / column)
    if n > most then n = math.ceil(w / tallest) end
    if n < 1 then n = 1 end
    return n
  end
}
\newcommand{\FrankVertPass}{%
  \par\addvspace{1.6ex}%
  \begingroup
    \FrankVertFont\SzPlain
    \setbox1\hbox{\foreignlanguage{japanese}{\FrankBuf}}%
    % the usable heights are the caps less the em of slack every column gets,
    % so a balanced column never exceeds its cap; \fontdimen6 is the em
    \dimen1=\dimexpr\FrankVertColumn-\fontdimen6\font\relax
    \dimen2=\dimexpr\FrankVertHeight-\fontdimen6\font\relax
    \count1=\directlua{tex.print(frank_vert_columns(\number\wd1, \number\dimen1,
      \number\dimen2, \number\linewidth, \number\baselineskip))}\relax
    \dimen3=\dimexpr\wd1/\count1+\fontdimen6\font\relax
    \setbox0\hbox{\rotatebox[origin=c]{-90}{%
      \parbox{\dimen3}{\FrankVertFont\SzPlain\color{ink}%
        \begin{otherlanguage}{japanese}\FrankBuf\par\end{otherlanguage}}}}%
    \ifdim\wd0>\linewidth
      \setbox0\hbox{\resizebox{\linewidth}{!}{\box0}}%
    \fi
    \setbox0\hbox{\raise\dp0\box0}%
    \noindent\hfill\box0\par
  \endgroup}

% ---- the "How to read this" page -----------------------------------------------
\newcommand{\FrankHowTo}{%
Every passage comes five times.

\smallskip
\textbf{First} the whole sentence with its reading written small above each
word — furigana, as a children's edition prints them — and nothing else.
Try to read it. Get as far as you can and do not worry about what you cannot
get; the point is the attempt, made before any help arrives.

\smallskip
\textbf{Second} only the sound of it: the reading of each phrase in kana, a
gap between one phrase and the next, and no kanji at all. Read it aloud. The
gaps are where the sentence divides.

\smallskip
\textbf{Third} the same sentence in those phrases, Japanese on the left, and
on the right four lines: the reading in kana, then the same in rōmaji, then
how to look each word up, then what it means. Now you find out how much you
had right.

\smallskip
\textbf{Fourth} the same passage again plain, without the furigana, the way
Japanese is printed for adults. You have read it three times, so it should
carry you.

\smallskip
\textbf{Fifth} the same again set vertically, top to bottom, the columns
running from right to left — the way most Japanese books are set. You already
know what it says, which is the only comfortable way to learn a new direction.

\smallskip
The reading over a word is the reading of the whole word, its kana ending
included, not of each character; the kana of a phrase is the reading of the
whole phrase. The rōmaji are Hepburn: long vowels carry a macron
(\textit{\=o}, \textit{\=u}),
the particles are written as they are said (\textit{wa}, \textit{o},
\textit{e}), and a doubled consonant is the small \textit{tsu}. Verbs are
given as dictionary form, \textit{-masu} stem, \textit{-te} form, meaning —
the \textit{-te} form is the one you cannot always guess. In brackets after
the meaning comes the verb's class, which makes every form not printed:
\textit{godan} verbs change the last vowel (\pw{書く}, \pw{書かない}),
\textit{ichidan} verbs drop the \textit{-ru} (\pw{食べる}, \pw{食べない}), and
\pw{する} and \pw{来る} are \textit{irregular}; a noun made a verb with
\pw{する} (\pw{勉強する}) is \textit{suru}. Then \textit{tr.} or
\textit{intr.} where it is known — it matters most for a pair that differs
only in that (\pw{開ける} you open a door, \pw{開く} it opens) — and
\textit{hon.} or \textit{hum.} for a word of respect or of modesty.
}
